\begin{resumo}


O ambiente universitário contemporâneo tem apresentado um aumento significativo nos casos relacionados a problemas de saúde entre estudantes, incluindo questões de saúde mental, sedentarismo, estresse acadêmico e hábitos prejudiciais ao bem-estar. Esses fatores impactam diretamente o desempenho estudantil e evidenciam a necessidade de mecanismos eficazes de monitoramento e intervenção. Diante disso, este trabalho tem como objetivo o desenvolvimento de um aplicativo \textit{mobile} voltado à coleta e análise de dados sobre o humor, hábitos e saúde geral dentro de instituições de ensino superior (IES). A proposta busca auxiliar o direcionamento de políticas institucionais de promoção à saúde, além de oferecer aos estudantes meios de buscar ajuda dentro da IES. A metodologia adotada envolve revisão bibliográfica e documental sobre o uso de tecnologias em saúde, análises de aplicativos semelhantes e aplicação de questionários à comunidade acadêmica, a fim de levantar os requisitos funcionais e não funcionais da aplicação. Sendo assim, na próxima fase, será iniciado o desenvolvimento do software, como implementação de funcionalidades, a fim de desenvolver um meio para auxiliar alunos na adoção de hábitos saudáveis e promover a saúde mental no meio acadêmico.



%  O resumo deve ressaltar o objetivo, o método, os resultados e as conclusões 
%  do documento. A ordem e a extensão
%  destes itens dependem do tipo de resumo (informativo ou indicativo) e do
%  tratamento que cada item recebe no documento original. O resumo deve ser
%  precedido da referência do documento, com exceção do resumo inserido no
%  próprio documento. (\ldots) As palavras-chave devem figurar logo abaixo do
%  resumo, antecedidas da expressão Palavras-chave:, separadas entre si por
%  ponto e finalizadas também por ponto. O texto pode conter no mínimo 150 e 
%  no máximo 500 palavras, é aconselhável que sejam utilizadas 200 palavras. 
%  E não se separa o texto do resumo em parágrafos.

%  - falar do PROBLEMA de saúde geral (não só mental mas evidenciar problema com hábitos, sedentarismo etc) dos universitários
%  - evidenciar o tema que é a utilização de um app voltado pra coletar esses dados que identificam os empecilhos dos universitários
%  - o objetivo de desenvolver um app que colete dados e trate pra direcionar políticas intervencionistas na faculdade ou, em situações mais graves, direcione ao atendimento individual
%  - metodologia de pesquisa e de desenvolvimento
%  - resultados não teremos mas dá pra adicionar uma conclusão final do tcc1
 

 \vspace{\onelineskip}
    
 \noindent
 \textbf{Palavras-chaves}: Aplicativo \textit{Mobile}, Hábitos Saudáveis, Saúde Mental.
\end{resumo}
